% Vector PDF atlas: choose a page, or use the per-figure SVG/PNG exports.

% Figure 1: Paper versus reproduction: all 16 recognition cases
\begin{figure}[t]
\centering
\includegraphics[page=1,width=\linewidth]{CSA_figure_atlas_20261030.pdf}
\caption{ORL: ten splits; MNIST: three; +/- is sample SD. FERET: one fixed identity partition. Gray CSA uses the observed-probe-selected 2026-10-30 recovery; Gabor CSA24* uses HOSVD/context. Other rows retain reference protocols. Crop/distance details are in REPORT; author split aggregation is unknown. Paper values: Xu et al., IEEE TCSVT (2008), Tables I/II/IV, doi:10.1109/TCSVT.2007.903317.}
\label{fig:all_recognition_cases}
\end{figure}

% Figure 2: ORL: six paper-dimension recognition cases
\begin{figure}[t]
\centering
\includegraphics[page=2,width=\linewidth]{CSA_figure_atlas_20261030.pdf}
\caption{Gray CSA11x2->16 uses the current probe-selected crop/distance configuration; Gabor CSA24* uses HOSVD/context. Other ORL rows retain shared references, so this is not one common input protocol. See REPORT for full settings; paper aggregation is unknown. Paper values: Xu et al., IEEE TCSVT (2008), Tables I/II/IV, doi:10.1109/TCSVT.2007.903317.}
\label{fig:orl_paper_comparison}
\end{figure}

% Figure 3: MNIST: both gallery sizes, both representations
\begin{figure}[t]
\centering
\includegraphics[page=3,width=\linewidth]{CSA_figure_atlas_20261030.pdf}
\caption{MNIST uses the original three-seed references. Best paper configurations and our repeated means have different reporting conventions. Paper values: Xu et al., IEEE TCSVT (2008), Tables I/II/IV, doi:10.1109/TCSVT.2007.903317.}
\label{fig:mnist_paper_comparison}
\end{figure}

% Figure 4: FERET: three training sizes, PCA and CSA controls
\begin{figure}[t]
\centering
\includegraphics[page=4,width=\linewidth]{CSA_figure_atlas_20261030.pdf}
\caption{These are fixed shared-cosine controls, not validation-selected winners. One fixed identity partition per q; no repeated-identity SD or significance claim. CSA does not beat the reproduced PCA at q=50. Paper values: Xu et al., IEEE TCSVT (2008), Tables I/II/IV, doi:10.1109/TCSVT.2007.903317.}
\label{fig:feret_paper_comparison}
\end{figure}

% Figure 5: Numerical gaps from the paper across all 16 cases
\begin{figure}[t]
\centering
\includegraphics[page=5,width=\linewidth]{CSA_figure_atlas_20261030.pdf}
\caption{Orange: current ORL CSA representatives, gray recovery81.84 and Gabor HOSVD/context88.75. Gray settings were selected using observed probes; see REPORT for crop/distance details. Other values retain reference protocols. Paper best-result aggregation is unknown. Paper values: Xu et al., IEEE TCSVT (2008), Tables I/II/IV, doi:10.1109/TCSVT.2007.903317.}
\label{fig:all_paper_gaps}
\end{figure}

% Figure 6: ORL recovery: every one of the ten observed splits
\begin{figure}[t]
\centering
\includegraphics[page=6,width=\linewidth]{CSA_figure_atlas_20261030.pdf}
\caption{Current CSA representative: HOSVD/context 88.75 +/- 2.89\%. Full identity 88.81 +/- 2.97\% is a separate start-shape variant. The dashed line refers to CSA, not an original PCA24 result. Observed-split candidate comparisons are exploratory.}
\label{fig:orl_ten_splits}
\end{figure}

% Figure 7: ORL: geometry and initialization interactions
\begin{figure}[t]
\centering
\includegraphics[page=7,width=\linewidth]{CSA_figure_atlas_20261030.pdf}
\caption{Geometry bundles change eye distance and height together; their individual effects are not isolated. Full-square initialization begins with larger mode dimensions than the final ranks. None of these means reaches 90.9\%.}
\label{fig:orl_geometry_initialization}
\end{figure}

% Figure 8: ORL fixed factorial study: sigma by context
\begin{figure}[t]
\centering
\includegraphics[page=8,width=\linewidth]{CSA_figure_atlas_20261030.pdf}
\caption{Whiskers are split SD. The effects are descriptive on repeated images. Context improves CSA24, but reducing sigma under context lowers CSA390; the settings are not jointly optimal for every method.}
\label{fig:orl_sigma_context_factorial}
\end{figure}

% Figure 9: ORL three fixed inputs: improvements and regressions
\begin{figure}[t]
\centering
\includegraphics[page=9,width=\linewidth]{CSA_figure_atlas_20261030.pdf}
\caption{120 outer fits; all negative branches retained. Context is a field-of-view extension and TT+HE is a photometric extension. This ten-split batch differs from the twenty-split factorial study.}
\label{fig:orl_three_fixed_inputs}
\end{figure}

% Figure 10: Gallery-only staged selection: most gain comes early
\begin{figure}[t]
\centering
\includegraphics[page=10,width=\linewidth]{CSA_figure_atlas_20261030.pdf}
\caption{Inputs include context/power extensions. CSA24 stages A->C add only about 0.03pp, and other controls can decline. Inner validation improvement does not guarantee outer improvement; this is not a blind new dataset.}
\label{fig:orl_gallery_selection_stages}
\end{figure}

% Figure 11: Crop-only bounded optimization: small gain, negative controls
\begin{figure}[t]
\centering
\includegraphics[page=11,width=\linewidth]{CSA_figure_atlas_20261030.pdf}
\caption{Dark: matched baseline; orange: validation-selected final setting. CSA24 improves only 0.16pp, while CSA390 and both PCA controls decline. Do not substitute this result for the context extension.}
\label{fig:orl_crop_only_optimization}
\end{figure}

% Figure 12: ORL sample-size extension with the same five probes
\begin{figure}[t]
\centering
\includegraphics[page=12,width=\linewidth]{CSA_figure_atlas_20261030.pdf}
\caption{Ten splits; sample SD whiskers. More fitting/gallery images change both representation estimation and NN coverage. G5/P5 is an extension, not the original G2/P8 experiment; no original-paper threshold is drawn.}
\label{fig:orl_sample_size}
\end{figure}

% Figure 13: MNIST bounded selection: preserve the negative result
\begin{figure}[t]
\centering
\includegraphics[page=13,width=\linewidth]{CSA_figure_atlas_20261030.pdf}
\caption{CSA declines by 0.77pp at G30 and 0.32pp at G40. These are ten-seed paired results, not changes from the three-seed main-reference means. All failed stages remain archived.}
\label{fig:mnist_selection_negative}
\end{figure}

% Figure 14: FERET: validation winners and fixed controls differ
\begin{figure}[t]
\centering
\includegraphics[page=14,width=\linewidth]{CSA_figure_atlas_20261030.pdf}
\caption{One fixed identity partition per q; no split SD. For q200 CSA, the fixed-cosine 95.38\% is higher than the actual validation-selected 93.77\%; both are retained.}
\label{fig:feret_selection_vs_fixed}
\end{figure}

% Figure 15: DATER supervised tensor extension: a negative result
\begin{figure}[t]
\centering
\includegraphics[page=15,width=\linewidth]{CSA_figure_atlas_20261030.pdf}
\caption{DATER: 20 new models; CSA/PCA controls reused. 19/20 DATER runs hit the20-iteration cap. Strict free-trajectory failure and numerical sensitivity are recorded; this result does not establish that DATER is generally ineffective.}
\label{fig:dater_matched_comparison}
\end{figure}

% Figure 16: Shrinkage ablation: compare both selected and matched variants
\begin{figure}[t]
\centering
\includegraphics[page=16,width=\linewidth]{CSA_figure_atlas_20261030.pdf}
\caption{Ten splits; SD whiskers. Ranks/input settings may differ between independently selected ordinary and shrinkage pipelines. Only the same-winner alpha0 comparison isolates removal of regularization in that selected configuration.}
\label{fig:shrinkage_ablation}
\end{figure}

% Figure 17: Representation cost: small weights, explicitly scoped timing
\begin{figure}[t]
\centering
\includegraphics[page=17,width=\linewidth]{CSA_figure_atlas_20261030.pdf}
\caption{Weight counts830 vs618792; arrays228272 vs5171968 bytes. Timing points are per-seed medians, not an uncertainty estimate. Excludes alignment, Gabor, training and I/O; counted arrays are not peak system memory.}
\label{fig:cost_scope}
\end{figure}

% Figure 18: Probe-eye perturbations: local robustness with no refitting
\begin{figure}[t]
\centering
\includegraphics[page=18,width=\linewidth]{CSA_figure_atlas_20261030.pdf}
\caption{Twenty overlapping splits. The full robustness batch has1280 scores and409600 repeated predictions, zero new fits. This does not cover single-eye errors, rotation or large offsets; no universal robustness claim.}
\label{fig:eye_perturbation}
\end{figure}

% Figure 19: Error analysis: corrections coexist with new failures
\begin{figure}[t]
\centering
\includegraphics[page=19,width=\linewidth]{CSA_figure_atlas_20261030.pdf}
\caption{3200 repeated probe records cover400 distinct photos. Input update:224 corrected and111 introduced errors. Same-budget comparison: CSA-only158, PCA-only249. Examples do not establish one causal explanation.}
\label{fig:error_tradeoffs}
\end{figure}

% Figure 20: Selected configurations: mean versus observed maximum
\begin{figure}[t]
\centering
\includegraphics[page=20,width=\linewidth]{CSA_figure_atlas_20261030.pdf}
\caption{Historical winning configurations were selected using probe scores across configurations and splits. Their means are also post-selection summaries. Different dimensions, inputs and split counts prevent a common-protocol ranking; maxima are not stable test estimates.}
\label{fig:mean_vs_maximum}
\end{figure}

% Figure 21: ORL tensor reconstruction: error and representation storage
\begin{figure}[t]
\centering
\includegraphics[page=21,width=\linewidth]{CSA_figure_atlas_20261030.pdf}
\caption{Historical Gabor/alignment version; not the new HOSVD/context recognition input. Learning-sample tensor reconstruction, not Gabor inversion or unseen-face generalization. Original absolute Gabor scale/pointwise curve values are unavailable; compare qualitative trends.}
\label{fig:orl_reconstruction}
\end{figure}

% Figure 22: CSA convergence: rapid stabilization in the recorded input
\begin{figure}[t]
\centering
\includegraphics[page=22,width=\linewidth]{CSA_figure_atlas_20261030.pdf}
\caption{This diagnostic does not replace the recognition limitTmax3. Rapid SSE/basis stabilization does not prove a global optimum or optimal classification. Exact paper ordinates and author input scaling are not recovered.}
\label{fig:csa_convergence}
\end{figure}

% Figure 23: Equal dimension does not imply equal retained variation
\begin{figure}[t]
\centering
\includegraphics[page=23,width=\linewidth]{CSA_figure_atlas_20261030.pdf}
\caption{These are projected-energy ratios, not percentages of identity information or accuracy. PCA and CSA have different structural constraints and centering conventions. This diagnostic alone does not identify the cause of the recognition gap.}
\label{fig:retained_variation}
\end{figure}

% Figure 24: FERET main experiment: recognition versus LDA dimension
\begin{figure}[t]
\centering
\includegraphics[page=24,width=\linewidth]{CSA_figure_atlas_20261030.pdf}
\caption{One fixed identity partition. These curves are descriptive evaluations across dimensions; do not reselect the winner using the held-out curve. They correspond to our fixed control, not pointwise recovery of original Fig.6.}
\label{fig:feret_lda_curves}
\end{figure}

% Figure 25: ORL historical dimension curves: three representation settings
\begin{figure}[t]
\centering
\includegraphics[page=25,width=\linewidth]{CSA_figure_atlas_20261030.pdf}
\caption{Historical alignment/Gabor/classification settings differ from the current88.75\% HOSVD/context representative. Do not splice the current endpoint into these early curves or label all curves as one final pipeline.}
\label{fig:orl_historical_dimension_curves}
\end{figure}

% Figure 26: CSA pipeline and the full reproduction scope
\begin{figure}[t]
\centering
\includegraphics[page=26,width=\linewidth]{CSA_figure_atlas_20261030.pdf}
\caption{Independent implementation of the paper method within the approved scope. Current ORL representative uses HOSVD initialization and context extension; full-square initialization is separate. Numerical agreement, implementation coverage and author-identical protocol are distinct claims.}
\label{fig:method_pipeline}
\end{figure}

% Figure 27: ORL current input: LDA dimension curves from saved models
\begin{figure}[t]
\centering
\includegraphics[page=27,width=\linewidth]{CSA_figure_atlas_20261030.pdf}
\caption{Descriptive scoring of saved directions1-23; no re-fitting or new parameter selection. Every23D endpoint matches the archived results. The fixed23D representative is retained even if a different curve point happens to score higher. Already-observed probes, context extension.}
\label{fig:orl_current_lda_curves}
\end{figure}
